\documentclass[main]{subfiles}

\begin{document}

\chapter{\MakeUppercase{EXPECTED OUTPUT}}

This chapter presents the expected outcomes of the proposed project on predicting urban growth and LULC dynamics using temporal machine learning models applied to Landsat (30 m) time-series data. The expected outputs are defined in terms of data products, analytical results, model performance, and practical applicability.

% \section{Historical LULC Maps}

% The project is expected to generate consistent multi-temporal land use/land cover maps for the selected study area using Landsat satellite imagery. These maps will represent major land-cover classes such as built-up areas, agricultural land, forest, water bodies, and barren land. The historical LULC maps will provide a clear visualization of urban expansion and land-cover transitions over time and serve as the foundation for temporal modeling.

\section{Temporal LULC Dataset}

A structured temporal dataset is expected to be produced by converting classified LULC maps into pixel-wise or patch-based time-series sequences. This dataset will represent the evolution of land-cover classes over multiple years and will be suitable for training and evaluating temporal machine learning models. The prepared dataset will be reusable for future studies and comparative experiments. Furthermore the dataset provides a cluster or densification zones which can help researchers for comparative studies and future research, providing a consistent foundation for analyzing urban growth dynamics in rapidly urbanizing regions such as the Kathmandu Valley.

\section{Trained Temporal Machine Learning Models}

The project is expected to develop trained temporal machine learning models, including LSTM and CNN--LSTM architectures, capable of learning spatio-temporal patterns from Landsat time-series data. These models will capture nonlinear temporal dependencies in land-cover transitions without relying on traditional CA--Markov assumptions.

\section{Future LULC and Urban Growth Predictions}

One of the key expected outputs is the prediction of future land use/land cover states over short- to medium-term horizons (e.g., 3--5 years). The predicted LULC maps will highlight potential urban expansion zones and areas susceptible to urban densification, including infill development and intensification within existing built-up areas. By capturing both horizontal expansion and internal urban restructuring, the predictions will provide a more comprehensive view of urban growth dynamics beyond peripheral sprawl.

\section{Model Performance Evaluation Results}

The project is expected to produce quantitative evaluation results that assess the predictive performance of temporal machine learning models. These results will include accuracy metrics such as overall accuracy, class-wise accuracy, confusion matrices, and Kappa statistics. Comparative analysis with baseline models, such as persistence and CA--Markov approaches, will demonstrate the relative strengths and limitations of temporal learning methods.

\section{Visualizations and Analytical Outputs}

The expected outputs also include visual products such as thematic LULC maps, change maps, and comparative plots illustrating historical trends and predicted future scenarios. Graphical representations of urban growth trends as well as urban densification of the existing build-up areas and model performance metrics will aid interpretation and communication of results.

\section{Practical and Research Contributions}

The project is expected to deliver a scalable and data-driven framework for urban growth and LULC prediction using freely available Landsat imagery and open-source machine learning tools. The methodology and outputs can support urban planners, policymakers, and researchers by providing predictive insights for land-use planning and sustainable urban development, particularly in data-scarce environments such as Nepal.

\end{document}
